\chapter{From a physical voltage to a sampled sequence}\label{ch:sampling}
\begin{readergoal}Understand the measurement chain, units, sample rates, aliasing, and the epistemic difference between source and signal.\end{readergoal}
\section{What electrodes actually measure}
Extracellular electrodes record potential differences relative to a reference. Their outputs reflect contributions of multiple transmembrane currents, electrode geometry, and conductive tissue. The recorded LFP is not a direct readout of a particular neuron or synapse. A physical forward model connects currents to voltages \cite{Buzsaki2012,Einevoll2013}. Analog amplification, referencing, and filtering occur before the digitized values are stored.

\section{Sampled time is a lattice}
For sampling frequency $f_s$, write $\Delta t=1/f_s$ and $x[n]=x(n\Delta t)$ for integer $n$ in the acquisition index set. The Nyquist frequency is $f_s/2$. A sinusoid at a frequency above Nyquist can produce the same discrete samples as a lower-frequency sinusoid; this is \emph{aliasing}. A pre-sampling low-pass filter is therefore required to suppress frequencies that would fold into the band of interest. Merely applying a digital low-pass filter after sampling cannot undo aliasing already introduced.

\section{The ideal sampling theorem and its assumptions}
If a continuous signal is strictly bandlimited to frequencies $|f|<f_s/2$, its samples determine the original signal through ideal sinc interpolation. Actual brain recordings are neither exactly bandlimited nor free of instrumentation artifacts. The theorem establishes a mathematical possibility under stated conditions, not a promise that every sampled waveform reproduces all physical current sources.

\section{Storage, quantization, and memory mapping}
A discrete recording is represented in finite precision. With $C$ channels, $N$ samples, and $b$ bytes per stored sample, uncompressed size is $C N b$ bytes (ignoring metadata and headers). Memory-mapped storage can support access to selected time blocks without allocating the entire dataset. It does not guarantee fast noncontiguous access: array layout, strides, page faults, and disk bandwidth still matter. Always document sample type, endian convention, orientation, sampling rate, units, and channel metadata.

\begin{workedexample}[A large acquisition]
For 128 channels at 30,000 samples/s over one hour, stored as signed 16-bit integers, the raw size is $128\cdot30000\cdot3600\cdot2=27{,}648{,}000{,}000$ bytes (about 27.65 GB in decimal units). Converting the entire array to 64-bit floats multiplies memory by four. A streaming analysis therefore need not sacrifice mathematical precision in its intermediate accumulators, but should avoid careless full-array expansion.
\end{workedexample}
\begin{exercises}
\item Determine the Nyquist frequency for an LFP sampled at 1250 Hz.
\item Calculate the storage size of 80 channels sampled at 2 kHz for 2 hours using float32.
\item Explain why the reference electrode can affect estimated coherence.
\item Distinguish a physical source model, an amplifier signal, and a floating-point array.
\end{exercises}